\section{设计约束与待定问题}

\subsection{约束与风险概述}
本系统当前处于课程项目迭代阶段，已经形成图像上传、图像检测、人工审核、通知和管理端治理的主要闭环，但论文检测、Review 检测、统一资源抽象和生产级部署仍需要后续完善。因此本章从当前约束、技术风险、业务风险、待定问题和后续工作五个角度，对概要设计中尚未完全落地的部分进行说明。

\subsection{当前约束}
\begin{itemize}[leftmargin=2em]
    \item 现有实现以学术图像检测为核心，论文与 Review 能力仍需扩展；
    \item 当前数据模型与接口设计中存在图像中心化特征，\texttt{ImageUpload}、\texttt{DetectionResult} 和人工审核链路仍是主要承载对象；
    \item AI 推理链路同时保留 HTTP 调用和旧版 SSH/SCP 调用方式，后续需要进一步收敛为稳定的服务间协议；
    \item 本地开发环境可使用 SQLite 和 eager 任务模式，但生产环境必须使用 MySQL、Redis 和独立 Worker，二者行为存在差异；
    \item 部分管理能力和模型配置能力尚未形成统一抽象，当前更偏向用户、组织、文件、日志和审核管理；
    \item 当前配置中仍存在开发态设置和硬编码敏感信息，生产化前必须改造为环境变量或私有配置。
\end{itemize}

\subsection{技术风险}
\begin{longtable}{|p{3.2cm}|p{6.1cm}|p{4.4cm}|}
\hline
风险项 & 风险说明 & 应对策略 \\ \hline
AI 服务稳定性 & 真实模型推理依赖 GPU、模型文件和复杂 Python 环境，存在启动失败、超时和返回格式不一致风险。 & 使用 HTTP 桩服务先稳定协议；真实模型接入时增加健康检查、超时、重试和结果 schema 校验。 \\ \hline
异步任务一致性 & Celery 任务失败可能导致任务状态、检测结果、报告文件和通知不一致。 & 增加失败状态、幂等写入和任务重试策略，关键步骤记录日志。 \\ \hline
WebSocket 部署 & Channels WebSocket 需要 ASGI 服务承载，若生产环境仅部署 WSGI，实时通知会失效。 & 明确 REST API 与 WebSocket 的运行进程和 Nginx 转发规则。 \\ \hline
文件存储增长 & 上传文件、结果图片和 PDF 报告持续增长，可能造成磁盘压力。 & 增加存储目录规划、归档策略、容量监控和历史任务清理策略。 \\ \hline
数据库迁移 & 当前模型偏图像场景，后续扩展统一资源模型时可能影响已有数据。 & 采用兼容迁移，优先新增表和映射关系，保留旧接口过渡期。 \\ \hline
配置泄露 & 数据库口令、邮箱授权码、服务 token 等若硬编码在仓库中，存在泄露风险。 & 迁移到环境变量或私有配置，并在部署文档中明确配置来源。 \\ \hline
\end{longtable}

\subsection{业务与管理风险}
\begin{itemize}[leftmargin=2em]
    \item \textbf{检测结论解释不足：} 若结果只给出真假判定，用户难以理解和复核，应保留置信度、可疑区域、子方法结果和文字说明；
    \item \textbf{人工审核边界不清：} 自动检测和人工审核结论可能不一致，需要明确展示两类结论的来源、时间和审核人；
    \item \textbf{组织权限误用：} 组织管理员、发布者和审核员权限边界复杂，若接口校验不充分，可能出现越权查看；
    \item \textbf{配额扣减争议：} 检测任务失败、重复提交或部分成功时，需要明确是否扣减配额以及如何补偿；
    \item \textbf{报告归档一致性：} 报告作为用户可下载凭证，必须与数据库结果一致，不能出现页面结果和 PDF 报告内容不一致；
    \item \textbf{模型版本追溯不足：} 若后续模型不断迭代，检测结果需要记录模型版本和规则版本，否则历史结果难以解释。
\end{itemize}

\subsection{待定问题}
\begin{itemize}[leftmargin=2em]
    \item 统一资源抽象与现有图像表结构的兼容迁移方案；
    \item 论文检测与 Review 检测的结果数据结构；
    \item 综合报告中多类型资源结果的整合方式；
    \item 模型管理配置在管理端中的可视化表达方式；
    \item 论文与 Review 专用日志的存储粒度和查询维度；
    \item 任务失败、部分成功、用户取消和管理员重跑等异常状态的完整状态机；
    \item AI 推理服务的正式鉴权方式、网络隔离方式和多模型路由方式；
    \item 文件归档、过期清理和对象存储迁移的具体策略；
    \item 生产环境中 WSGI/ASGI、Celery Worker、Redis 和 Nginx 的最终部署组合。
\end{itemize}

\subsection{后续迭代建议}
\begin{enumerate}[leftmargin=2em]
    \item \textbf{第一阶段：收敛当前图像检测闭环。} 完善任务失败状态、报告生成状态、AI 返回 schema 校验和部署配置环境变量化，保证当前主流程稳定；
    \item \textbf{第二阶段：抽象统一资源与统一任务。} 在不破坏现有图像接口的前提下，引入资源类型、任务类型和差异化结果结构，为论文与 Review 检测接入做准备；
    \item \textbf{第三阶段：接入论文检测与 Review 检测。} 优先完成文本解析、段落级结果、风险解释和报告聚合，再逐步接入真实模型；
    \item \textbf{第四阶段：完善管理端治理。} 增加模型配置、模型版本、任务失败原因、资源归档和组织级统计能力；
    \item \textbf{第五阶段：生产化部署加固。} 完成 Nginx、ASGI、Celery、Redis、MySQL、AI 服务和文件存储的正式部署脚本、监控告警和备份恢复流程。
\end{enumerate}

\subsection{后续工作}
后续将基于本概要设计继续细化详细设计，包括数据库设计、接口字段设计、模块内部类与函数设计、关键状态机、接口时序图、异常处理策略、测试设计和部署脚本。具体工作包括：
\begin{itemize}[leftmargin=2em]
    \item 给出统一资源、统一检测任务、论文结果、Review 结果和综合报告的详细数据表设计；
    \item 补齐接口字段、请求示例、响应示例、错误码和权限说明；
    \item 细化 Celery 任务拆分、队列命名、重试策略和失败恢复策略；
    \item 明确 AI 推理服务协议，包括健康检查、批量检测、模型版本、返回 schema 和错误响应；
    \item 完善本地、集成、生产三类环境的部署文档和启动脚本；
    \item 设计覆盖上传、检测、审核、通知、管理统计和报告下载的测试用例。
\end{itemize}

\subsection{本章小结}
本章对当前概要设计中的约束、风险和待定问题进行了集中说明。总体来看，项目现有图像检测链路已经具备继续扩展的基础，但后续需要重点解决统一资源抽象、AI 服务协议收敛、生产部署加固、任务异常恢复和多类型检测报告整合等问题。通过分阶段推进，可以在保持当前功能可用的同时，逐步演进到更完整的学术 AI 鉴伪平台。
